% !TeX root = ../../Book1.tex
% 纲要标题：## 第15章　分裂域、代数闭包与正规扩张
% 纲要位置：计划纲要.md，第 528 行。

\chapter{分裂域、代数闭包与\mbox{正规扩张}}
\label{b1:ch:15}
\BookChapterNumber{b1:ch:15}

\begin{chapterabstract}
分裂域把一个多项式的全部根置于同一最小扩张中，代数闭包则容纳所有代数方程的根。本章构造这两类域并证明相应的唯一性，研究嵌入延拓，最后用共轭根刻画正规扩张。
\end{chapterabstract}

\begin{learninggoals}
\begin{itemize}
\item 构造并计算低次多项式的分裂域，区分同构类型唯一与同构映射唯一。
\item 解释代数闭包的存在性构造，以及 Zorn 引理在极大理想和部分嵌入中的作用。
\item 用分裂条件、嵌入条件和根的生成条件判断正规性。
\item 构造正规闭包，把共轭根之间的同构延拓到更大的代数扩张。
\end{itemize}
\end{learninggoals}

多项式 \(X^3-2\) 有实根 \(\sqrt[3]{2}\)，但域 \(\mathbb Q(\sqrt[3]{2})\) 仍没有它的另外两个复根。若只为方程找到一个根，便无法同时研究三个根之间的代数关系；继续添入三次单位根后，才得到包含全部根的最小扩张。不同的添根次序又可能给出外观不同的域。分裂域的构造与唯一性要回答的，正是如何在不依赖添根次序的情况下讨论方程的所有根。

\ProseParagraphBreak

本章使用\cref{b1:prop:minimal-polynomial-degree}的单扩张结构及\cref{b1:thm:algebraic-tower}的代数性传递。分裂域只需有限次添根；代数闭包与一般嵌入延拓使用 Zorn 引理。实数和复数仅用于具体低次算例；分裂域的存在性证明适用于任意特征。

\ProseParagraphBreak
分裂域、嵌入和正规性可结合 Lang 的《Algebra》\cite[Chapter V, \S\S 2--4]{Lang2002}阅读；Roman 的《Field Theory》\cite[Sections 2.7--2.9, 3.1--3.2]{Roman2006FieldTheory}也可用于比较逐次添根与嵌入延拓。比较不同教材的证明时，应特别核对它们处理的是有限扩张，还是任意代数扩张，以及何处需要选择原理。

% !TeX root = ../../Book1.tex
% 纲要标题：### 15.1 分裂域
% 纲要位置：计划纲要.md，第 530 行。

\section{分裂域}
\label{b1:sec:1501}


% 纲要内容（仅作写作计划，尚非教材正文）：
%
% * 多项式的分裂域构造
% * 嵌入延拓与分裂域的唯一性；系数变换与根的像须保持已有关系
% * 具体低次多项式的分裂域
%

添入一个根使某个不可约因子变为一次因子，却未必得到其他根。例如实数 \(\sqrt[3]2\) 不会使 \(x^3-2\) 在实数域中完全分裂。若要同时研究全部根之间的对称性，应选取由全部根生成的域；它避免了任意附加与原方程无关的元素。

\subsection{多项式的分裂域构造}
\label{b1:subsec:1501:01}

\begin{definition}\label{b1:def:splitting-field}
设 \(K\) 为域，\(n\geq1\) 为整数，\(f\in K[x]\) 为 \(n\) 次多项式。如果域扩张 \(L/K\) 中存在 \(a\in K^\times\) 及 \(\alpha_1,\ldots,\alpha_n\in L\)，使
\[
f=a\prod_{i=1}^n(x-\alpha_i),\qquad L=K(\alpha_1,\ldots,\alpha_n),
\]
则称 \(L\) 为 \(f\) 在 \(K\) 上的\term[fenlie-yu]{分裂域}[splitting field]。非零常数多项式的分裂域约定为 \(K\)。
\end{definition}
第一个条件要求全部根都已进入 \(L\)，第二个条件要求 \(L\) 恰由这些根生成。例如 \(X^2-2\) 在 \(\mathbb C\) 中分裂，但 \(\mathbb C\) 不是它在 \(\mathbb Q\) 上的分裂域；两个根只生成 \(\mathbb Q(\sqrt2)\)。根按重数列出，重复列出一个根并不改变生成的域。

\begin{proposition}\label{b1:prop:splitting-field-exists}
设 \(K\) 为域。每个非零多项式 \(f\in K[x]\) 都有分裂域，且分裂域为 \(K\) 的有限扩张。若 \(\deg f=n\geq1\)，则该扩张的次数至多为 \(n!\)。
\end{proposition}
\begin{proof}
若 \(f\) 非常数，取一个不可约因子 \(q\)。域 \(K_1=K[x]/(q)\) 包含元素 \(\alpha_1=x+(q)\)，它满足 \(q(\alpha_1)=0\)，从而也是 \(f\) 的根。这里商环为域可由多项式 Bézout 关系直接验证：非零剩余类的代表元与 \(q\) 互素，故有逆元。由因式定理，在 \(K_1[x]\) 中写 \(f=(x-\alpha_1)f_1\)。

在 \(K_1\) 上重复此构造，为次数 \(n-1\) 的 \(f_1\) 添入一个根。每除去一个一次因子，剩余多项式的次数就减少一，所以至多 \(n\) 步便得到由全部这些根生成、使 \(f\) 分裂的域。每步所选不可约因子的次数不超过当时剩余多项式的次数，故各步扩张次数分别至多为 \(n,n-1,\ldots,1\)。由\cref{b1:thm:tower-law}知总次数至多为 \(n!\)。若某个根已经在当前域内，就可继续除去相应的一次因子而不扩大域；重根只需按其重数继续除去，不要求反复添入同一个元素。
\end{proof}
这是一个存在性构造。实际计算不必机械地逐根相添；根之间已知的关系常能显著减少步骤，也能把 \(n!\) 这个粗上界缩小。

\subsection{嵌入延拓与分裂域的唯一性}
\label{b1:subsec:1501:02}

若将基域元素按 \(\sigma\) 变换，那么关系 \(m(\alpha)=0\) 的系数也随之变换。因此 \(\alpha\) 的像必须满足 \(\sigma(m)(\beta)=0\)，不能仅凭 \(m(\beta)=0\) 选择它。下面的引理说明，这一条件也足以得到单扩张上的延拓。

\begin{lemma}\label{b1:lem:simple-embedding-extension}
设 \(L/K\) 为域扩张，\(\alpha\in L\) 在 \(K\) 上代数，最小多项式为 \(m\in K[x]\)。设 \(\sigma:K\rightarrow K'\) 为域同构，\(\Omega/K'\) 为域扩张，并记 \(\sigma(m)\in K'[x]\) 为对 \(m\) 的系数施行 \(\sigma\) 所得的多项式。若预先选定 \(\sigma(m)\) 在 \(\Omega\) 中的一个根 \(\beta\)，则存在唯一的域嵌入 \(\widetilde\sigma_\beta:K(\alpha)\rightarrow\Omega\)，使
\[
\widetilde\sigma_\beta|_K=\sigma,\qquad \widetilde\sigma_\beta(\alpha)=\beta.
\]
反过来，\(\sigma\) 的每个此类延拓都把 \(\alpha\) 送到 \(\sigma(m)\) 的一个根。因此延拓与 \(\sigma(m)\) 在 \(\Omega\) 中的不同根一一对应。
\end{lemma}
\begin{proof}
若 \(\tau:K(\alpha)\rightarrow\Omega\) 延拓 \(\sigma\)，令 \(\beta=\tau(\alpha)\)。把 \(m(\alpha)=0\) 经 \(\tau\) 映过去即得 \(\sigma(m)(\beta)=0\)。反过来，给定这样的根 \(\beta\)，系数同构把不可约 \(m\) 变为不可约 \(\sigma(m)\)。由\cref{b1:prop:minimal-polynomial-degree}，\(K(\alpha)=K[\alpha]\)，故可定义
\[
\widetilde\sigma_\beta:\sum_i a_i\alpha^i\longmapsto\sum_i\sigma(a_i)\beta^i.
\]
若两个多项式表达式在 \(\alpha\) 处相等，其差 \(h\in K[x]\) 属于代入映射的核 \((m)\)，即 \(h=qm\)。变换系数后有 \(\sigma(h)=\sigma(q)\sigma(m)\)，在 \(\beta\) 处仍为零，所以像与表达式的选择无关。该映射保持加法、乘法和单位，并在域之间单射；其像为 \(K'(\beta)\)。每个元素都是 \(\alpha\) 的多项式，故在已选定 \(\beta\) 后，延拓只能是 \(\widetilde\sigma_\beta\)。
\end{proof}

例如取 \(K=\mathbb Q(\sqrt2)\)、\(\alpha=\sqrt[4]2>0\)，并在 \(\mathbb C\) 中计算。由\cref{b1:thm:eisenstein}，\(X^4-2\) 在 \(\mathbb Q\) 上不可约，故 \([\mathbb Q(\alpha):\mathbb Q]=4\)。又有 \(\sqrt2=\alpha^2\)，塔式公式给出 \([K(\alpha):K]=2\)，所以 \(\alpha\) 在 \(K\) 上的最小多项式为
\[
m=X^2-\sqrt2.
\]
令 \(\sigma:K\rightarrow K\) 把 \(\sqrt2\) 送到 \(-\sqrt2\)，则
\[
\sigma(m)=X^2+\sqrt2,
\qquad \sigma(m)\;\text{的根为}\;\pm i\alpha.
\]
因此 \(\sigma\) 从 \(K(\alpha)\) 到 \(\mathbb C\) 的两个嵌入延拓必须分别把 \(\alpha\) 送到 \(i\alpha\) 与 \(-i\alpha\)。若送到 \(\pm\alpha\)，平方仍为 \(\sqrt2\)，就不满足系数已经变换后的关系。这里定义域为实域 \(K(\alpha)=\mathbb Q(\alpha)\)，目标域为 \(\mathbb C\)，两个延拓的像域均为 \(K(i\alpha)\)，含有非实元素。嵌入的定义域、目标域与像域是不同的对象。

\ProseParagraphBreak
即使基域保持不变，也不能独立指定各生成元的共轭根。例如 \(\alpha\) 与 \(b=\alpha^2\) 在 \(\mathbb Q\) 上分别满足 \(X^4-2\) 与 \(X^2-2\)。若一个 \(\mathbb Q\) 同态固定 \(\alpha\)，便必须把 \(b\) 送到 \(\alpha^2=b\)，不能另选 \(-b\)。逐次延拓时要使用当前中间域上的最小多项式：在 \(\mathbb Q(\alpha)\) 上，\(b\) 的最小多项式已是 \(X-\alpha^2\)，没有另选符号的自由。

\begin{proposition}\label{b1:prop:splitting-field-unique}
设 \(f\in K[x]\) 为非零多项式，\(L/K\) 是 \(f\) 的分裂域，\(\sigma:K\rightarrow K'\) 是域同构，而 \(L'/K'\) 是 \(\sigma(f)\) 的分裂域。则 \(\sigma\) 可以延拓为域同构 \(L\rightarrow L'\)。特别地，同一多项式的两个分裂域在保持基域的同构意义下唯一，但这样的同构映射未必唯一。
\end{proposition}
\begin{proof}
先逐根延拓基域同构以取得嵌入，再利用目标域恰由全部根生成来证明它满射。把 \(L\) 写成 \(K(\alpha_1,\ldots,\alpha_r)\)，其中所列元素是 \(f\) 的全部不同根。假设已在 \(E=K(\alpha_1,\ldots,\alpha_{i-1})\) 上构造嵌入 \(\tau:E\rightarrow L'\)。因为 \(f(\alpha_i)=0\)，\cref{b1:prop:minimal-polynomial-degree}说明 \(\alpha_i\) 在 \(E\) 上的最小多项式整除 \(f\)。所以 \(\tau(m_{\alpha_i,E})\) 整除 \(\sigma(f)\)，后者在 \(L'\) 中分裂。于是前者在 \(L'\) 中有根，由\cref{b1:lem:simple-embedding-extension}可把 \(\tau\) 延拓到 \(E(\alpha_i)\)。有限次延拓给出嵌入 \(L\rightarrow L'\)。

在 \(L[x]\) 中，\(f\) 的一次分解经此嵌入变为 \(\sigma(f)\) 的一次分解。因此 \(\sigma(f)\) 的全部根都在像域中，像域又包含 \(K'\)。\(L'\) 由这些根生成，故像恰为 \(L'\)，嵌入为同构。
\end{proof}
唯一的是同构类型，不是同构本身。例如 \(\mathbb Q(\sqrt2)\) 到自身有保持 \(\mathbb Q\) 的两个同构，分别把 \(\sqrt2\) 送到 \(\sqrt2\) 与 \(-\sqrt2\)。

\subsection{具体低次多项式的分裂域}
\label{b1:subsec:1501:03}

在 \(\mathbb Q\) 上，\(x^2-2\) 的两个根是 \(\pm\sqrt2\)，分裂域为 \(\mathbb Q(\sqrt2)\)，次数二。对 \(x^4-2\)，令 \(a=\sqrt[4]2>0\)，其根为 \(a,-a,ia,-ia\)，故分裂域为 \(\mathbb Q(a,i)\)。Eisenstein 判据给出 \(\bigl[\mathbb Q(a):\mathbb Q\bigr]=4\)，而 \(\mathbb Q(a)\subseteq\mathbb R\) 不含 \(i\)，所以再添 \(i\) 的次数为二，总次数为八。

\ProseParagraphBreak
对 \(x^3-2\)，令 \(a=\sqrt[3]2>0\)，\(\zeta=(-1+\sqrt{-3})/2\)。其三个根为 \(a,a\zeta,a\zeta^2\)，故分裂域为 \(\mathbb Q(a,\zeta)\)。\(x^3-2\) 在 \(\mathbb Q\) 上不可约，而 \(\zeta\) 不实且满足 \(x^2+x+1=0\)，所以
\begin{equation}\label{b1:eq:cubic-splitting-degree}
\begin{aligned}
\bigl[\mathbb Q(a,\zeta):\mathbb Q\bigr]
&=\bigl[\mathbb Q(a,\zeta):\mathbb Q(a)\bigr]\,\bigl[\mathbb Q(a):\mathbb Q\bigr]\\
&=2\cdot3=6.
\end{aligned}
\end{equation}
根 \(a\) 生成的三次扩张并未包含另外两个根；分裂域把它们一起纳入，才为后续置换全部根提供合适的域。添一个根的次数由它的最小多项式控制，分裂域的次数还反映其余根需要补充多少代数信息。

\ProseParagraphBreak
由单扩张嵌入引理，\(a\mapsto a\zeta\) 给出 \(\mathbb Q(a)\) 到 \(\mathbb C\) 的 \(\mathbb Q\) 嵌入，其像为 \(\mathbb Q(a\zeta)\)。这个像含有非实元素，而 \(\mathbb Q(a)\) 是实域，所以两个子域不同，尽管它们保持 \(\mathbb Q\) 地同构。同构描述代数结构相同，并不意味着在共同环境中占据同一个子域。

% !TeX root = ../../Book1.tex
% 纲要标题：### 15.2 代数闭包
% 纲要位置：计划纲要.md，第 536 行。

\section{代数闭包}
\label{b1:sec:1502}


% 纲要内容（仅作写作计划，尚非教材正文）：
%
% * 代数闭包的存在；有限多条根关系的相容性与递增并域
% * 一般嵌入延拓与代数闭包在基域同构意义下的唯一性
% * 选择原理在证明中的位置
%

一个分裂域只处理指定多项式。若要在同一个域中讨论基域上的所有代数方程，可先构造代数扩张，使 \(K[x]\) 中的每个非常数多项式至少有一个根。但得到新域 \(F_1\) 后，\(F_1[x]\) 中还可能出现系数不全属于 \(K\) 的多项式，刚才的构造并未为这些方程安排根。因此对新域继续添根，得到递增的域列 \(F_0=K\subseteq F_1\subseteq F_2\subseteq\cdots\)。利用每个多项式只有有限多个系数，可以证明这些域的并处理了其中全部的多项式方程。

\subsection{代数闭包的存在性}
\label{b1:subsec:1502:01}

\begin{definition}\label{b1:def:algebraic-closure}
设 \(\overline K/K\) 为域扩张。如果 \(\overline K/K\) 是代数扩张，并且 \(\overline K\) 是代数闭域，即 \(\overline K[x]\) 中每个非常数多项式都在 \(\overline K\) 中有根，则称 \(\overline K\) 为 \(K\) 的\term[daishu-bibao]{代数闭包}[algebraic closure]。
\end{definition}
\symidx{\(\overline K\)：\(K\) 的一个代数闭包}

要使域 \(F\) 上的每个非常数多项式都有根，可以为每个首一非常数多项式 \(f\) 设置一个形式符号 \(X_f\)，再在多项式环中施加关系 \(f(X_f)=0\)。若这些关系在一个非零商环中同时成立，\(X_f\) 的像便成为 \(f\) 的根。下面还要使这个商环成为域，并保证所得扩张代数。

\begin{lemma}\label{b1:lem:adjoin-all-roots}
对任意域 \(F\)，存在代数扩张 \(F'/F\)，使 \(F[x]\) 中每个非常数多项式在 \(F'\) 中至少有一个根。
\end{lemma}
\begin{proof}
为每个首一非常数多项式 \(f\in F[x]\) 设置一个不定元 \(X_f\)，令 \(R=F[X_f\mid f]\)。每个环元素只涉及有限多个不定元。令 \(I\) 为全部 \(f(X_f)\) 生成的理想；在 \(R/I\) 中，这些关系便同时成立。

我们先证 \(I\neq R\)。关键是任何有限多个这样的关系都能在某个扩域中同时实现。如果 \(1\in I\)，按生成理想的定义，就有有限表达式
\[
1=\sum_{i=1}^r h_i f_i(X_{f_i}),\qquad h_i\in R,
\]
其中可将重复的 \(f_i\) 合并。由\cref{b1:prop:splitting-field-exists}逐次为 \(f_1,\ldots,f_r\) 添入根，得到某个扩张域 \(E/F\)，在其中同时选定 \(f_i\) 的根 \(a_i\)。把 \(X_{f_i}\) 送到 \(a_i\)，其余不定元送到零，就得到保持 \(F\) 的环同态 \(R\rightarrow E\)。余变量取零只是为了给整个多项式环定义同态，并不要求这次赋值满足它们各自对应的添根关系；上式只需所列的有限条关系成立。上式的右边映为零，左边仍为一，矛盾。因此有限子系统的相容性排除了 \(1\in I\)。

由\cref{b1:prop:maximal-ideal-exists}，真理想 \(I\) 包含于一个极大理想 \(\mathfrak m\)。令 \(F'=R/\mathfrak m\)，则它是域；\(F\cap\mathfrak m=0\)，因为非零常数可逆，所以 \(F\) 嵌入 \(F'\)。记 \(X_f\) 的像为 \(\overline X_f\)，则 \(f(\overline X_f)=0\)，所以每个 \(\overline X_f\) 都在 \(F\) 上代数。任意 \(a\in F'\) 都是有限多个 \(\overline X_{f_1},\ldots,\overline X_{f_s}\) 的多项式表达，因而属于
\[
F(\overline X_{f_1},\ldots,\overline X_{f_s}).
\]
由\cref{b1:lem:finite-algebraic-generators}，这个中间域在 \(F\) 上有限，再由\cref{b1:lem:finite-dimensional-algebraic}知 \(a\) 在 \(F\) 上代数。因此 \(F'/F\) 是代数扩张。任意非常数多项式除以首项系数便成为首一多项式，根不改变，故所有非常数多项式都有根。
\end{proof}

\begin{theorem}\label{b1:thm:algebraic-closure}
每个域 \(K\) 都存在代数闭包。
\end{theorem}
\begin{proof}
从 \(F_0=K\) 出发，反复应用\cref{b1:lem:adjoin-all-roots}构造代数扩张 \(F_{n+1}/F_n\)，使 \(F_n\) 上每个非常数多项式在 \(F_{n+1}\) 中有根。将各步嵌入认同为包含，令 \(\Omega=\bigcup_{n\geq0}F_n\)。任意两个元素都落在某个共同的 \(F_n\) 中，其差、积以及非零元素的逆仍在该域内，因此 \(\Omega\) 是域。反复应用\cref{b1:thm:algebraic-tower}可知每个 \(F_n/K\) 代数。任意 \(a\in\Omega\) 属于某个 \(F_n\)，所以在 \(K\) 上代数；这证明了 \(\Omega/K\) 代数。

若 \(g=\sum_{i=0}^d c_ix^i\in\Omega[x]\) 非常数，每个系数 \(c_i\) 属于某个 \(F_{n_i}\)。系数只有有限多个，取 \(n=\max\{n_0,\ldots,n_d\}\)，并利用域列递增，就有所有 \(c_i\in F_n\)。所以 \(g\in F_n[x]\)，在 \(F_{n+1}\subseteq\Omega\) 中有根。故 \(\Omega\) 代数闭，满足代数闭包的两项要求。
\end{proof}
这项构造不要求基域可数。每一步同时为该域的全部多项式安排不定元；迭代后的有限系数论证直接证明并域代数闭，不需要先判定第一步所得域是否已代数闭。“可数次迭代”也不意味着所得代数闭包是可数集。

\subsection{一般嵌入延拓与代数闭包的唯一性}
\label{b1:subsec:1502:02}

\begin{theorem}[嵌入延拓]\label{b1:thm:embedding-extension}
设 \(L/K\) 为任意代数扩张，\(\Omega\) 为代数闭域，\(\sigma:K\rightarrow\Omega\) 为域嵌入。则存在嵌入 \(\widetilde\sigma:L\rightarrow\Omega\) 延拓 \(\sigma\)。
\end{theorem}
\begin{proof}
只要一个部分嵌入的定义域还没有覆盖 \(L\)，就尝试再添入一个元素并延拓。Zorn 引理给出极大部分嵌入，而代数性将保证其定义域不能遗漏任何元素。

令 \(\mathcal P\) 为所有 \((E,\tau)\) 组成的集合，其中 \(E\) 是 \(L/K\) 的中间域，\(\tau:E\hookrightarrow\Omega\) 是满足 \(\tau|_K=\sigma\) 的域嵌入。规定 \((E,\tau)\le(E',\tau')\) 当且仅当 \(E\subseteq E'\) 且 \(\tau'|_E=\tau\)。初始对 \((K,\sigma)\) 使 \(\mathcal P\) 非空。

对一条非空链 \(\mathcal C\)，将其定义域取并，得到中间域 \(E_{\mathcal C}\)。链中任意两个定义域按包含关系可比，较大域上的映射限制到较小域后就是原映射，所以两个映射在重叠处绝不冲突。于是可定义 \(\tau_{\mathcal C}(a)=\tau(a)\)，其中任选一个包含 \(a\) 的链成员 \((E,\tau)\)，所得值与选择无关。任意有限多个元素都可放进同一个链成员的定义域，在那里验证 \(\tau_{\mathcal C}\) 保持域运算并且单射，就知它是延拓 \(\sigma\) 的嵌入。故 \((E_{\mathcal C},\tau_{\mathcal C})\) 给出上界；空链以 \((K,\sigma)\) 为上界。Zorn 引理（\cref{b1:thm:A-zorn}）给出极大对 \((E,\tau)\)。

若 \(E\neq L\)，取 \(a\in L\setminus E\)。因为 \(L/K\) 代数，\(a\) 满足某个非零的 \(K\) 系数多项式方程。\(K\subseteq E\)，所以同一个方程也是非零的 \(E\) 系数方程，\(a\) 因而在 \(E\) 上代数。其最小多项式经系数同构 \(\tau:E\rightarrow\tau(E)\) 变为 \(\tau(E)\) 上的不可约多项式。\(\Omega\) 代数闭，故该多项式在 \(\Omega\) 中有根。\cref{b1:lem:simple-embedding-extension}便把 \(\tau\) 延拓到 \(E(a)\)，与极大性矛盾。所以 \(E=L\)，得到所需延拓。
\end{proof}

\begin{proposition}\label{b1:prop:algebraic-closure-unique}
若 \(\Omega/K\)、\(\Omega'/K\) 都是代数闭包，则存在保持 \(K\) 的同构 \(\Omega\rightarrow\Omega'\)。
\end{proposition}
\begin{proof}
对代数扩张 \(\Omega/K\) 和嵌入 \(K\hookrightarrow\Omega'\) 应用\cref{b1:thm:embedding-extension}，得到保持 \(K\) 的嵌入 \(\tau:\Omega\rightarrow\Omega'\)。为了证明像域 \(\tau(\Omega)\) 代数闭，任取其中的非常数多项式
\[
g(x)=\sum_{i=0}^d\tau(a_i)x^i,
\qquad a_i\in\Omega,
\]
将系数拉回，得到 \(g^\tau(x)=\sum_{i=0}^d a_ix^i\in\Omega[x]\)。因 \(\tau\) 单射，这个多项式仍非常数。\(\Omega\) 代数闭，故有 \(c\in\Omega\) 使 \(g^\tau(c)=0\)。映回去便有
\[
g(\tau(c))=\tau\bigl(g^\tau(c)\bigr)=0,
\]
所以 \(g\) 在像域中有根。

任意 \(b\in\Omega'\) 在 \(K\) 上代数。由于 \(\tau\) 保持 \(K\)，有 \(K\subseteq\tau(\Omega)\)；原来零化 \(b\) 的非零 \(K\) 系数多项式仍是非零的 \(\tau(\Omega)\) 系数多项式，所以扩大基域以后 \(b\) 仍代数。其在像域上的不可约最小多项式只能为一次，因为像域代数闭。于是 \(b\in\tau(\Omega)\)，所以 \(\tau\) 满射。这里利用了“代数闭域没有非平凡代数扩张”，而非错误地认为无限集上的单射自动为满射。
\end{proof}
同构通常并不唯一；选择不同共轭根会给出不同延拓。后文写“固定一个代数闭包”时，是选定共同环境，后续的嵌入与根都在这个选定的域中讨论。

\subsection{证明中的选择原理}
\label{b1:subsec:1502:03}

\cref{b1:lem:adjoin-all-roots}由真理想取得极大理想，调用了先前用 Zorn 引理证明的极大理想存在性；\cref{b1:thm:embedding-extension}则直接对部分嵌入使用 Zorn 引理。

\ProseParagraphBreak
对链取并时必须检查两件不同的事：定义域之并对域运算封闭；映射在重叠处给出同一个值。只有在偏序要求“延拓”而不只是“定义域更大”时，第二件事才得到保证。极大部分嵌入的论证还依赖 \(L/K\) 代数：若剩余元素超越，不能再以最小多项式及其根作为延拓依据。

\ProseParagraphBreak
本书在这些一般存在性论证中采用通常的选择公理。算法性问题另作处理：抽象地得到一个极大理想或代数闭包，并未给出可计算的元素编码、多项式分解程序或根的选择算法。具体基域上的有效计算，须另加适当表示方法。

% !TeX root = ../../Book1.tex
% 纲要标题：### 15.3 正规扩张
% 纲要位置：计划纲要.md，第 542 行。

\section{正规扩张}
\label{b1:sec:1503}


% 纲要内容（仅作写作计划，尚非教材正文）：
%
% * 分裂条件与嵌入条件；代数扩张到自身的嵌入自动为自同构
% * 正规闭包；有限生成元最小多项式的分裂域与全部嵌入像的生成域
% * 正规性对中间域和塔的行为
%

分裂域包含一个方程的所有根。正规扩张把这一要求推广为：凡是不可约方程已有一个根落入扩张，其余根也必须一并落入。用嵌入来表述时，这意味着改变根的选择不会把整个扩张带到另一个子域。本节的扩张均为代数扩张。

\ProseParagraphBreak
给定代数扩张 \(L/K\)，由\cref{b1:thm:algebraic-closure}可取 \(L\) 的代数闭包 \(\Omega\)。因为 \(\Omega/L\) 与 \(L/K\) 都代数，\cref{b1:thm:algebraic-tower}说明 \(\Omega/K\) 代数；\(\Omega\) 又是代数闭域，所以它同时是包含指定 \(L\) 的 \(K\) 的代数闭包。下文的根与嵌入均在这样选定的 \(\Omega\) 中讨论。

\subsection{分裂条件与嵌入条件}
\label{b1:subsec:1503:01}

\begin{definition}\label{b1:def:normal-extension}
设 \(L/K\) 为代数扩张。如果每个在 \(L\) 中有根的 \(K\) 系数不可约多项式都在 \(L\) 中分裂，则称 \(L/K\) 为\term[zhenggui-kuozhang]{正规扩张}[normal extension]。
\end{definition}
这里的分裂允许重根：只要求多项式成为一次因子的乘积，不要求这些因子互不相同。正规性关心所有根是否留在域内；不同根的个数是下一章可分性所研究的问题。

\begin{proposition}\label{b1:prop:algebraic-self-embedding}
设 \(L/K\) 为代数扩张，\(\sigma:L\rightarrow L\) 为逐点固定 \(K\) 的域嵌入。则 \(\sigma\) 为 \(L\) 的 \(K\) 自同构，无须假设 \(L/K\) 正规。
\end{proposition}
\begin{proof}
即使 \(L/K\) 无限，每个元素仍受一个有限次数的方程约束。给定 \(a\in L\)，令 \(R\) 为 \(m_{a,K}\) 在 \(L\) 中的不同根组成的集合。它有限且包含 \(a\)；\(\sigma\) 固定多项式的系数，故将 \(R\) 单射地映入自身。有限集上的单射必为双射，所以存在 \(b\in R\) 使 \(\sigma(b)=a\)。每个 \(a\in L\) 都有原像，故 \(\sigma\) 满射。
\end{proof}
本命题的目标域已经是 \(L\)，所以只须解决满射性。正规性的嵌入判据还须保证：当目标扩大为代数闭包时，像不会离开 \(L\)。

\begin{theorem}\label{b1:thm:normal-embedding-criterion}
设 \(L/K\) 为代数扩张，\(\Omega\) 为包含 \(L\) 的 \(K\) 的代数闭包。以下条件等价：
\begin{enumerate}
\item \(L/K\) 正规。
\item 每个保持 \(K\) 的嵌入 \(\sigma:L\rightarrow\Omega\) 都满足 \(\sigma(L)=L\)。
\item 存在一族 \(K[x]\) 中的非零多项式，使 \(L\) 在 \(K\) 上由它们在 \(\Omega\) 中的全部根生成。
\end{enumerate}
有限扩张时，这些条件还等价于 \(L\) 是某一个非零 \(K\) 系数多项式的分裂域。
\end{theorem}
\begin{proof}
先设正规。对任意 \(a\in L\)，其最小多项式在 \(L\) 中分裂，\(\sigma(a)\) 也是该多项式的根，所以 \(\sigma(L)\subseteq L\)。将目标限制为 \(L\) 后，\cref{b1:prop:algebraic-self-embedding}给出满射性，故 \(\sigma(L)=L\)。这里适用的是逐元素的有限根集论证，既不要求 \(L/K\) 有限，也没有把一般无限集上的单射当作满射。

设第二个条件成立。若不可约 \(f\in K[x]\) 在 \(L\) 中有根 \(a\)，要证明它分裂，就须把任意另一个根 \(b\in\Omega\) 也纳入 \(L\)。\cref{b1:lem:simple-embedding-extension}给出 \(K(a)\rightarrow K(b)\)、\(a\mapsto b\) 的 \(K\) 同构，再由\cref{b1:thm:embedding-extension}延拓为 \(L\rightarrow\Omega\)。其像按假设为 \(L\)，故 \(b\in L\)。因此全部根属于 \(L\)，证明正规性。

正规时，取 \(L\) 中所有元素在 \(K\) 上的最小多项式；这些非零多项式的全部根都在 \(L\) 中，并在 \(K\) 上生成 \(L\)，得到第三个条件。反过来，若 \(L\) 在 \(K\) 上由一族非零多项式在 \(\Omega\) 中的全部根生成，则每个 \(K\) 嵌入把每个有限根集单射地映入自身，因而置换该根集。生成域遂被映到自身且映满，满足第二个条件。

若 \(L/K\) 有限，取有限生成元 \(a_1,\ldots,a_r\)，令 \(f=\prod_{i=1}^r m_{a_i,K}\)。正规性使各最小多项式都在 \(L\) 中分裂，故全部根生成的域包含于 \(L\)；这个域又包含各 \(a_i\)，因而包含 \(L\)。所以 \(L\) 正是 \(f\) 的分裂域；反方向已由第三个条件得到。
\end{proof}
第三个条件必须取所选多项式的全部根。仅由其中一些根生成的域可以是任意代数扩张，未必正规。例如 \(\mathbb Q(\sqrt2)/\mathbb Q\) 正规；\(\mathbb Q(\sqrt[3]2)/\mathbb Q\) 不正规，因为 \(x^3-2\) 的另两个根不实。此处三次扩张和六次分裂域之间的差别，正是正规性所检测的差别。

\subsection{正规闭包}
\label{b1:subsec:1503:02}

\begin{proposition}\label{b1:prop:normal-closure}
设 \(L/K\) 为代数扩张，\(\Omega\) 为包含 \(L\) 的 \(K\) 的代数闭包。则在所选 \(\Omega\) 中，按子域的包含关系存在包含 \(L\) 的最小正规扩张 \(N/K\)。它由所有 \(a\in L\) 的最小多项式在 \(\Omega\) 中的全部根生成，称为 \(L/K\) 在 \(\Omega\) 中的\term[zhenggui-bibao]{正规闭包}[normal closure]。它也有描述
\[
N=K\left(\bigcup_{\sigma\in\operatorname{Hom}_K(L,\Omega)}\sigma(L)\right),
\]
其中 \(\operatorname{Hom}_K(L,\Omega)\) 表示逐点固定 \(K\) 的嵌入 \(L\rightarrow\Omega\) 全体。若 \(L/K\) 有限且 \(L=K(a_1,\ldots,a_r)\)，则 \(N\) 正是 \(\prod_{i=1}^r m_{a_i,K}\) 在 \(\Omega\) 中的分裂域，特别地 \(N/K\) 有限。
\end{proposition}
\begin{proof}
按所述根集生成 \(N\)，\cref{b1:thm:normal-embedding-criterion}的第三个条件直接保证 \(N/K\) 正规，且每个 \(a\in L\) 都在生成根集中，故 \(L\subseteq N\)。若正规扩张 \(M/K\)、\(L\subseteq M\subseteq\Omega\)，则每个 \(a\) 的最小多项式在 \(M\) 中分裂，全部生成根都属于 \(M\)，故 \(N\subseteq M\)。

记所有像域生成的域为 \(M_0\)。每个 \(K\) 嵌入将 \(a\in L\) 送到 \(m_{a,K}\) 的一个根，故 \(\sigma(L)\subseteq N\)，从而 \(M_0\subseteq N\)。反过来，任取 \(m_{a,K}\) 在 \(\Omega\) 中的根 \(\beta\)。\cref{b1:lem:simple-embedding-extension}给出 \(a\mapsto\beta\) 的 \(K(a)\rightarrow\Omega\) 嵌入，再由\cref{b1:thm:embedding-extension}延拓到 \(L\)。于是 \(\beta\) 属于某个像域，故 \(N\) 的全部生成根都属于 \(M_0\)。这证明 \(N=M_0\)：正规闭包把 \(L\) 在共同代数闭包中的所有 \(K\) 同构像一起纳入。

若 \(L=K(a_1,\ldots,a_r)\) 为有限扩张，取 \(m_{a_i,K}\) 的乘积在 \(\Omega\) 内的分裂域 \(N_0\)。它有限、正规并包含 \(L\)，故 \(N\subseteq N_0\)。另一方面，\(N/K\) 正规且包含各 \(a_i\)，所以包含各 \(m_{a_i,K}\) 的全部根，故 \(N_0\subseteq N\)。因此 \(N=N_0\) 为有限扩张；这个具体描述不依赖生成集的选择。
\end{proof}
例如 \(\mathbb Q(\sqrt[3]2)\) 的正规闭包就是 \(\mathbb Q(\sqrt[3]2,\zeta)\)，其中 \(\zeta^2+\zeta+1=0\)、\(\zeta\neq1\)。这里先固定共同代数闭包，才可以谈“最小子域”；若不固定环境，结论应改写成保持基域和已给扩张的同构意义下唯一。

\subsection{正规性与中间域、扩张塔}
\label{b1:subsec:1503:03}

\begin{proposition}\label{b1:prop:normal-tower-behavior}
设 \(K\subseteq E\subseteq L\) 为代数扩张塔。若 \(L/K\) 正规，则 \(L/E\) 正规。在同一代数闭包内，任意一族正规 \(K\) 扩张的交以及它们生成的合成域都对 \(K\) 正规。但 \(L/K\) 正规不保证 \(E/K\) 正规，\(E/K\) 与 \(L/E\) 都正规也不保证 \(L/K\) 正规。
\end{proposition}
\begin{proof}
每个 \(E\) 嵌入 \(L\rightarrow\Omega\) 也是 \(K\) 嵌入，故由\cref{b1:thm:normal-embedding-criterion}保持 \(L\)，从而 \(L/E\) 正规。

若 \(a\) 属于一族正规域之交，其在 \(K\) 上的最小多项式在每个域中分裂，所以全部根属于交域，证明交域正规。对于合成域，每个被合成的域由某族 \(K\) 系数多项式的根生成；把各族合并即可满足\cref{b1:thm:normal-embedding-criterion}的第三个条件。

第一个反例取 \(L=\mathbb Q(\sqrt[3]2,\zeta)\)、\(E=\mathbb Q(\sqrt[3]2)\)、\(K=\mathbb Q\)，其正规与非正规已在本节说明。第二个取
\[
K=\mathbb Q,\quad E=\mathbb Q(\sqrt2),\quad L=\mathbb Q(\sqrt[4]2).
\]
\(E\) 是 \(x^2-2\) 的分裂域，而 \(L/E\) 是 \(x^2-\sqrt2\) 的分裂域，故两层正规。但 \(x^4-2\) 在 \(\mathbb Q\) 上不可约，其非实根不在实域 \(L\) 中，故总扩张不正规。
\end{proof}
交和合成是在同一环境中组合若干扩张，域塔则是在一个扩张内改变基域；前两种运算保持正规性，并不意味着正规性沿塔传递。将来用子群描述中间域时，这些不同的行为会分别对应不同的群论条件。

% !TeX root = ../../Book1.tex
% 纲要标题：### 15.4 嵌入的延拓与共轭
% 纲要位置：计划纲要.md，第 548 行。

\section{嵌入的延拓与共轭}
\label{b1:sec:1504}


% 纲要内容（仅作写作计划，尚非教材正文）：
%
% * 中间域嵌入的延拓与自同构；代数闭包上的自同构延拓
% * 多项式根的共轭性
% * 为自同构计数准备工具
%
% ---
%

比较两个根时，我们通常不只关心它们满足同一个方程，还关心替换这个根能否同时替换域中的其他元素。嵌入延拓定理保证代数关系可相容地延续；正规性决定延拓后的像是否仍是原域。本节把这两种作用分清，并整理随后计数所需的形式。

\subsection{中间域嵌入的延拓与自同构}
\label{b1:subsec:1504:01}

\begin{corollary}\label{b1:cor:intermediate-embedding-extension}
设 \(K\subseteq E\subseteq L\) 为域扩张塔，并且 \(L/K\) 代数；设 \(\Omega\) 为包含 \(L\) 的 \(K\) 的代数闭包。每个 \(K\) 嵌入 \(\tau:E\rightarrow\Omega\) 均能延拓为 \(K\) 嵌入 \(L\rightarrow\Omega\)。若 \(L/K\) 正规，将该延拓的目标限制为 \(L\) 后，它为 \(L\) 的 \(K\) 自同构。此外，\(\tau\) 总能延拓为 \(\Omega\) 的 \(K\) 自同构。
\end{corollary}
\begin{proof}
\(L/E\) 代数，对它与嵌入 \(\tau\) 使用\cref{b1:thm:embedding-extension}，便得到 \(\widetilde\tau:L\rightarrow\Omega\)。由于它延拓 \(\tau\)，在 \(K\) 上仍为恒等；当 \(L/K\) 正规时，\cref{b1:thm:normal-embedding-criterion}给出 \(\widetilde\tau(L)=L\)。所以将目标限制为 \(L\) 后，才得到 \(L\) 的自同构。

\(\Omega/K\) 代数，故 \(\Omega/E\) 也代数。再次应用嵌入延拓定理，可将 \(\tau\) 延拓为 \(\Omega\rightarrow\Omega\) 的 \(K\) 嵌入。\cref{b1:prop:algebraic-self-embedding}保证它满射，因而是 \(\Omega\) 的 \(K\) 自同构。
\end{proof}
上述结论也允许 \(E/K\) 无限；在代数闭包中，局部选定的共轭根可以相容地延拓到整个域的自同构。

\subsection{多项式根的共轭性}
\label{b1:subsec:1504:02}

\begin{definition}\label{b1:def:field-conjugate-elements}
设 \(\Omega\) 为域 \(K\) 的代数闭包，\(a,b\in\Omega\) 都在 \(K\) 上代数。如果 \(a,b\) 具有相同的 \(K\) 上最小多项式，则称它们为 \(K\) 上的\term[gonge-yuan]{共轭元}[conjugate elements]，也说它们在 \(K\) 上共轭。
\end{definition}
这等价于 \(a,b\) 是同一个 \(K[x]\) 中首一不可约多项式的根。这里的共轭是域论术语，其定义与群中的 \(gxg^{-1}\) 不同；下面通过域的自同构解释它。

\begin{proposition}\label{b1:prop:conjugates-automorphism}
设 \(\Omega\) 为域 \(K\) 的代数闭包，\(a,b\in\Omega\) 都在 \(K\) 上代数。则 \(a,b\) 在 \(K\) 上共轭，当且仅当存在 \(\Omega\) 的 \(K\) 自同构把 \(a\) 送到 \(b\)。若还有中间域 \(K\subseteq L\subseteq\Omega\) 包含 \(a,b\)，且 \(L/K\) 正规，则也存在 \(L\) 的 \(K\) 自同构实现这一替换。
\end{proposition}
\begin{proof}
同一个不可约多项式的两个根由\cref{b1:lem:simple-embedding-extension}给出 \(K(a)\rightarrow K(b)\) 的 \(K\) 同构，先在单扩张上实现所需替换。再应用\cref{b1:cor:intermediate-embedding-extension}，把它延拓为 \(\Omega\) 的 \(K\) 自同构；若 \(L/K\) 正规，也可将它延拓为 \(L\) 的 \(K\) 自同构。反过来，若自同构送 \(a\) 到 \(b\)，则把 \(m_{a,K}(a)=0\) 映过去，得到 \(m_{a,K}(b)=0\)。由于它首一不可约，必为 \(b\) 的最小多项式。
\end{proof}
共轭性依赖基域。例如 \(\sqrt2\) 与 \(-\sqrt2\) 在 \(\mathbb Q\) 上的最小多项式都是 \(X^2-2\)，但在 \(K=\mathbb Q(\sqrt2)\) 上分别为 \(X-\sqrt2\) 和 \(X+\sqrt2\)，不再相同。扩大基域后，可用的系数增加，原来不能区分的根可能被区分。即使在正特征中出现重根，命题谈的也是不同的元素；重数不会产生新的嵌入。

\subsection{嵌入个数与扩张次数}
\label{b1:subsec:1504:03}

保持 \(K\) 的域同态 \(L\rightarrow\Omega\) 全体记为 \(\operatorname{Hom}_K(L,\Omega)\)。域同态必单射，所以此处它就是嵌入集合。保持 \(K\) 的 \(L\) 的自同构全体记为 \(\operatorname{Aut}_K(L)\)，复合使它成为群。
\symidx{\(\operatorname{Hom}_K(L,\Omega)\)：保持 \(K\) 的域嵌入集合}
\symidx{\(\operatorname{Aut}_K(L)\)：保持 \(K\) 的域自同构群}

\begin{proposition}\label{b1:prop:embedding-degree-bound}
设 \(L/K\) 为有限域扩张，\(\Omega\) 为包含 \(L\) 的 \(K\) 的代数闭包。则
\[
1\leq\#\operatorname{Hom}_K(L,\Omega)\leq[L:K].
\]
若 \(L/K\) 还正规，则这些嵌入恰为 \(\operatorname{Aut}_K(L)\) 的元素。
\end{proposition}
\begin{proof}
\cref{b1:thm:embedding-extension}将 \(\operatorname{id}_K\) 延拓到 \(L\)，所以至少有一个嵌入。为证明上界，写 \(L=K(a_1,\ldots,a_r)\)，令 \(E_0=K\)、\(E_i=K(a_1,\ldots,a_i)\)。给定 \(E_{i-1}\) 的一个嵌入，它到 \(E_i\) 的延拓由 \(a_i\) 的最小多项式经系数变换后的不同根决定。该多项式的次数为 \([E_i:E_{i-1}]\)，故\cref{b1:lem:simple-embedding-extension}给出的延拓数至多为这个数。每个 \(E_i\) 的嵌入都限制为一个 \(E_{i-1}\) 的嵌入，不同限制所对应的延拓彼此不同，因此
\[
\#\operatorname{Hom}_K(E_i,\Omega)
\leq [E_i:E_{i-1}]\,\#\operatorname{Hom}_K(E_{i-1},\Omega).
\]
从 \(E_0\) 的唯一嵌入出发，逐层相乘并用\cref{b1:thm:tower-law}，便得 \(\#\operatorname{Hom}_K(L,\Omega)\leq[L:K]\)。正规情形的断言正是\cref{b1:thm:normal-embedding-criterion}：每个嵌入的像都等于 \(L\)。
\end{proof}
这个上界可能严格。取素数 \(p\)，令 \(t\) 在 \(\mathbb F_p\) 上超越，并置
\[
K=\mathbb F_p(t^p),\qquad L=\mathbb F_p(t).
\]
取包含 \(L\) 的 \(K\) 的代数闭包 \(\Omega\)。由\eqref{b1:eq:rational-power-degree}，\([L:K]=p\)，所以 \(t\) 在 \(K\) 上的最小多项式是 \(X^p-t^p\)。在特征 \(p\) 下，
\[
X^p-t^p=(X-t)^p.
\]
它在 \(\Omega\) 中只有一个不同的根 \(t\)。每个 \(K\) 嵌入 \(L\rightarrow\Omega\) 都须把 \(t\) 送到这个根，因而只能是恒等映射。于是
\[
\#\operatorname{Hom}_K(L,\Omega)=|\operatorname{Aut}_K(L)|=1<p=[L:K].
\]
同时，\(L\) 是 \(X^p-t^p\) 在 \(K\) 上的分裂域，故 \(L/K\) 正规。这个例子表明，正规性保证上述嵌入的像仍为 \(L\)，却不保证嵌入数达到扩张次数；重根的重数不会产生新的嵌入。下一章将证明：对有限扩张，如果没有不可分现象，嵌入数便等于扩张次数；若同时正规，这个数就是自同构群的阶。


\begin{samepage}
本章各类嵌入结论的条件列于\cref{b1:tab:embedding-normality-count}；达到次数上界的条件将在下一章讨论。
\begin{table}[H]
\centering
\small
\caption{嵌入延拓、正规性与计数的条件比较}\label{b1:tab:embedding-normality-count}
\begin{tabularx}{\linewidth}{@{}p{0.27\linewidth}X@{}}
\toprule
所处理的对象 & 假设与结论 \\
\midrule
单扩张 \(K(\alpha)/K\) & 给定域同构 \(\sigma:K\rightarrow K'\) 及扩张域 \(\Omega/K'\)，并选定 \(\sigma(m_{\alpha,K})\) 在 \(\Omega\) 中的根 \(\beta\) 后，\cref{b1:lem:simple-embedding-extension}给出唯一满足 \(\alpha\mapsto\beta\) 的延拓。 \\
一般代数扩张 \(L/K\) & 目标域 \(\Omega\) 代数闭时，给定的基域嵌入 \(K\hookrightarrow\Omega\) 可延拓为 \(L\hookrightarrow\Omega\)；延拓通常不唯一（\cref{b1:thm:embedding-extension}）。 \\
正规扩张 \(L/K\) & 在包含 \(L\) 的 \(K\) 的代数闭包 \(\Omega\) 中，每个逐点固定 \(K\) 的嵌入 \(\tau:L\hookrightarrow\Omega\) 都满足 \(\tau(L)=L\)（\cref{b1:thm:normal-embedding-criterion}）。 \\
有限扩张 \(L/K\) & 若 \([L:K]<\infty\)，则 \(L\) 到包含它的 \(K\) 的代数闭包 \(\Omega\) 中的 \(K\) 嵌入数至多为 \([L:K]\)（\cref{b1:prop:embedding-degree-bound}）。 \\
\bottomrule
\end{tabularx}
\end{table}
\end{samepage}

\begin{chaptersummary}
添入一个代数元时，嵌入的延拓由最小多项式的一个根决定。逐次使用这一事实，可构造分裂域并证明其在基域同构意义下唯一。把部分嵌入按延拓排序，则能将同样的想法推广到任意代数扩张；其中的极大性论证必须同时保证定义域的封闭性和映射的相容性。

\ProseParagraphBreak
设 \(L/K\) 为代数扩张，\(\Omega\) 为包含 \(L\) 的 \(K\) 的代数闭包。正规性可由一族非零 \(K\) 系数多项式的全部根生成来刻画；当 \(L/K\) 正规时，每个逐点固定 \(K\) 的嵌入 \(\sigma:L\hookrightarrow\Omega\) 都满足 \(\sigma(L)=L\)。这个判据也适用于无限代数扩张。正规性可以沿塔向上层限制，也对交和合成保持，但既不自动向中间域下降，也不沿两层正规扩张传递。

\ProseParagraphBreak
若 \([L:K]<\infty\)，则逐点固定 \(K\) 的嵌入 \(L\hookrightarrow\Omega\) 的数目至多为 \([L:K]\)；何时达到这个上界，将由下一章的可分性理论说明。
\end{chaptersummary}

\BookExercises

\begin{exercise}\label{b1:ex:15-1}
求 \(x^4+4\) 在 \(\mathbb Q\) 上的分裂域及其次数。
\end{exercise}
\begin{solution}
在 \(\mathbb C\) 中分解为 \((x^2-2x+2)(x^2+2x+2)\)，四个根为 \(1\pm i,-1\pm i\)。它们都属于 \(\mathbb Q(i)\)，而其中的根 \(1+i\) 减去 \(1\) 就得到 \(i\)，所以它们在 \(\mathbb C\) 中生成的子域恰为 \(\mathbb Q(i)\)。因此 \(\mathbb Q(i)\) 是一个分裂域；任意其他分裂域都与它保持 \(\mathbb Q\) 地同构。\(x^2+1\) 在 \(\mathbb Q\) 上无根，故次数为二。
\end{solution}

\begin{exercise}\label{b1:ex:15-2}
设 \(a=\sqrt[4]3>0\)，\(b=a^2=\sqrt3\)。
\begin{enumerate}
\item 求 \(x^4-3\) 在 \(\mathbb Q\) 上的分裂域，并列出 \(\mathbb Q(a)\) 到 \(\mathbb C\) 的全部 \(\mathbb Q\) 嵌入。
\item 分别求 \(a,b\) 在 \(\mathbb Q\) 上的最小多项式。虽然 \(a\mapsto a\)、\(b\mapsto-b\) 分别把它们送到各自最小多项式的根，证明这两个指定不能同时由一个 \(\mathbb Q\) 域同态实现，并说明在 \(\mathbb Q(a)\) 上为何不能再独立选择 \(b\) 的像。
\end{enumerate}
\end{exercise}
\begin{solution}
1.\ 根为 \(a,-a,ia,-ia\)，分裂域为 \(\mathbb Q(a,i)\)。\(x^4-3\) 对素数三满足 Eisenstein 判据，故 \(\bigl[\mathbb Q(a):\mathbb Q\bigr]=4\)。该域实，添入 \(i\) 再给出二次扩张，所以分裂域次数为八。由\cref{b1:lem:simple-embedding-extension}，\(\mathbb Q(a)\) 的嵌入恰由把 \(a\) 送到上述四个根给出；其中只有像为 \(a,-a\) 的两个嵌入把 \(\mathbb Q(a)\) 映入自身。

2.\ \(a,b\) 在 \(\mathbb Q\) 上的最小多项式分别为 \(X^4-3\) 与 \(X^2-3\)。所指定的两个像各自都是对应多项式的根，但若同态固定 \(a\)，就必须把 \(b=a^2\) 送到 \(a^2=b\)，不能送到 \(-b\)。在当前中间域 \(\mathbb Q(a)\) 上，\(b\) 的最小多项式是 \(X-a^2\)，所以逐次延拓时也没有另选其符号的自由。
\end{solution}

\begin{exercise}\label{b1:ex:15-3}
设 \(f\in K[x]\) 非零，\(r\geq1\) 为整数。证明 \(f\) 与 \(f^r\) 在同一个 \(K\) 的代数闭包中的分裂域相同。再取非零 \(g\in K[x]\)，并在同一个代数闭包中取 \(f,g\) 的分裂域；说明乘积 \(fg\) 的分裂域如何由它们得到。
\end{exercise}
\begin{solution}
\(f\) 与 \(f^r\) 有相同的不同根，只是重数改变，故由根生成的域相同。若 \(f,g\) 的分裂域分别为 \(L_f,L_g\)，乘积 \(fg\) 的根集是两个根集的并，故其分裂域是包含 \(L_f,L_g\) 的最小子域，即合成域 \(L_fL_g\)。这也包括常数多项式的情形。
\end{solution}

\begin{exercise}\label{b1:ex:15-4}
在 \(\mathbb F_2[x]\) 中考虑 \(f=x^2+x+1\)。不用有限域分类，直接构造其分裂域，并写出保持 \(\mathbb F_2\) 的两个自同构。
\end{exercise}
\begin{solution}
\(f(0)=f(1)=1\)，所以二次多项式不可约，\(L=\mathbb F_2[x]/(f)\) 是域。记 \(a=x+(f)\)，则 \(a^2+a+1=0\)，又 \(f(a+1)=f(a)=0\)，故它的两个根为 \(a,a+1\)，均在 \(L\) 中。由 \(L=\mathbb F_2(a)\) 知它就是分裂域。两个自同构分别为 \(a\mapsto a\) 与 \(a\mapsto a+1\)；后一映射重复两次恢复 \(a\)，也可直接验证为自同构。
\end{solution}

\begin{exercise}\label{b1:ex:15-5}
证明每个代数闭域都是无限域。若 \(K\) 可数，证明其任意代数闭包也可数。
\end{exercise}
\begin{solution}
若域 \(F\) 有限，则 \(1+\prod_{a\in F}(x-a)\) 为非常数多项式，在每个 \(a\in F\) 处值为一，故没有根，\(F\) 不是代数闭域。若 \(K\) 可数，\(K[x]\) 中非零多项式只有可数多个，每个在代数闭包中只有有限多个不同根，而代数闭包的每个元素都属于某个这样的有限根集。因此代数闭包至多可数；由刚证的无限性，它恰为可数无限集。
\end{solution}

\begin{exercise}\label{b1:ex:15-6}
设 \(\Omega\) 是代数闭域，\(K\subseteq\Omega\)，不要求 \(\Omega/K\) 代数。证明 \(\Omega\) 中在 \(K\) 上代数的元素组成 \(K\) 的一个代数闭包。
\end{exercise}
\begin{solution}
记该集合为 \(A\)。\cref{b1:prop:algebraic-elements-subfield}说明它是子域，且 \(A/K\) 按定义代数。给定非常数 \(f\in A[x]\)，在 \(\Omega\) 中取根 \(a\)。其有限多个系数在 \(K\) 上生成有限扩张 \(E\)，而 \(a\) 在 \(E\) 上代数。\cref{b1:thm:algebraic-tower}给出 \(a\) 在 \(K\) 上代数，故 \(a\in A\)。于是 \(A\) 代数闭，同时满足代数闭包的两个条件。
\end{solution}

\begin{exercise}\label{b1:ex:15-7}
设 \(L/K\) 代数，\(\Omega\) 为 \(K\) 的代数闭包。证明每个 \(K\) 嵌入 \(L\rightarrow\Omega\) 的像都是一个与 \(L\) 同构的中间域，但这些像未必相同。给出具体例子。
\end{exercise}
\begin{solution}
域嵌入保持运算与逆元，故像是包含 \(K\) 的子域，映射本身给出到像的 \(K\) 同构。取 \(L=\mathbb Q(a)\)、\(a=\sqrt[3]2>0\)，并在一个包含三个根的代数闭包中取非实根 \(a\zeta\)。嵌入 \(a\mapsto a\zeta\) 的像为 \(\mathbb Q(a\zeta)\)，包含非实元素；原来的 \(L\) 全部实，故二者不相同。它们尽管嵌入在同一环境中位置不同，抽象的 \(\mathbb Q\) 域结构仍同构。
\end{solution}

\begin{exercise}\label{b1:ex:15-8}
设 \(L/K\) 正规，\(E/K\) 是同一代数闭包中的任意代数扩张。证明合成域 \(LE/E\) 正规。若 \(L/K\) 有限，证明 \([LE:E]\leq[L:K]\)。
\end{exercise}
\begin{solution}
\(L\) 由某族 \(K\) 系数多项式的全部根生成，故 \(LE\) 由 \(E\) 与同一族根生成；把系数视为 \(E\) 中元素，\cref{b1:thm:normal-embedding-criterion}即说明 \(LE/E\) 正规。有限情形下，取 \(L/K\) 的基 \(u_1,\ldots,u_n\)。它们的 \(E\) 线性生成空间包含 \(E,L\)，且对乘法封闭，因为 \(u_iu_j\) 可用 \(K\) 系数写成该基的线性组合。非零元素在这个有限维空间上的乘法为单射，因而满射，所以它也有逆元。因此该空间是域，恰等于 \(LE\)，从而 \([LE:E]\leq n\)。
\end{solution}

\begin{exercise}\label{b1:ex:15-9}
证明每个二次扩张都正规，不论其特征如何。
\end{exercise}
\begin{solution}
设 \([L:K]=2\)，取 \(a\in L\setminus K\)。塔式公式给出 \(L=K(a)\)，其最小多项式写成 \(x^2+bx+c\)。已知一个根 \(a\)，另一个根为 \(-b-a\)，也在 \(L\) 中，包括二者相等的情形。因此 \(L\) 是该多项式的分裂域，由\cref{b1:thm:normal-embedding-criterion}正规。证明没有除以二，所以也适用于特征二。
\end{solution}

\begin{exercise}\label{b1:ex:15-10}
设 \(\sigma:L\rightarrow L\) 为保持 \(K\) 的域嵌入，\(L/K\) 代数，但未必正规。证明 \(\sigma\) 必为自同构。给出删去代数性后的反例。
\end{exercise}
\begin{solution}
给定 \(a\in L\)，令 \(R\) 为 \(m_{a,K}\) 在 \(L\) 中的不同根组成的有限非空集合。\(\sigma\) 固定系数，因此将 \(R\) 单射地映入自身，必为双射，故 \(a\) 有 \(R\) 中原像。对每个 \(a\) 都如此，所以 \(\sigma\) 满射。反例是 \(k(t)\rightarrow k(t)\)、\(t\mapsto t^2\)：超越性保证这是嵌入，但像为 \(k(t^2)\)，由\eqref{b1:eq:rational-power-degree}知它是次数二的真子域。
\end{solution}

\begin{exercise}\label{b1:ex:15-11}
设 \(E/K\) 代数，\(\Omega/K\) 为代数闭包且 \(E\subseteq\Omega\)。证明 \(\Omega\) 也是 \(E\) 的代数闭包，并说明 \(E\) 的任意 \(K\) 嵌入到 \(\Omega\) 能否延拓为 \(\Omega\) 的自同构。
\end{exercise}
\begin{solution}
\(\Omega\) 本来代数闭；其中元素在 \(K\) 上的代数方程也可视为 \(E\) 系数方程，所以 \(\Omega/E\) 代数。因此它是 \(E\) 的代数闭包。给定嵌入，用\cref{b1:thm:embedding-extension}延拓到 \(\Omega\rightarrow\Omega\)，所得映射仍固定 \(K\)。它是代数扩张 \(\Omega/K\) 到自身的嵌入，由上一题的有限根集论证必满射，故为自同构；也可直接用\cref{b1:cor:intermediate-embedding-extension}之后的结论。
\end{solution}

\begin{exercise}\label{b1:ex:15-12}
令 \(a=\sqrt[3]2>0\)，\(\zeta^2+\zeta+1=0\)、\(\zeta\neq1\)，\(L=\mathbb Q(a,\zeta)\)。证明复共轭在 \(L\) 上给出自同构，并求它的固定元素。
\end{exercise}
\begin{solution}
共轭固定 \(a\)，并把 \(\zeta\) 送到 \(\zeta^2=-1-\zeta\)，故保持生成域 \(L\)，其平方为恒等。\eqref{b1:eq:cubic-splitting-degree}给出 \(\bigl[L:\mathbb Q(a)\bigr]=2\)，所以元素唯一写成 \(u+v\zeta\)、\(u,v\in\mathbb Q(a)\)。其共轭为 \(u-v-v\zeta\)，与原式相等便由系数比较得 \(v=-v\)，故 \(v=0\)。所以固定元素恰为 \(\mathbb Q(a)\)。
\end{solution}

\begin{exercise}\label{b1:ex:15-13}
设 \(L/K\) 为有限扩张，整数 \(r\geq1\)，且 \(L=K(a_1,\ldots,a_r)\)。固定包含 \(L\) 的 \(K\) 的代数闭包 \(\Omega\)，并在其中记 \(L\) 的正规闭包为 \(N\)。证明 \(N\) 是 \(\prod_{i=1}^r m_{a_i,K}\) 在 \(\Omega\) 中的分裂域。进一步证明
\[
N=K\left(\bigcup_{\sigma\in\operatorname{Hom}_K(L,\Omega)}\sigma(L)\right),
\]
并由此说明 \(N\) 不依赖生成集的选择。
\end{exercise}
\begin{solution}
令 \(N_0\) 为题中多项式在 \(\Omega\) 中的分裂域。它正规并包含全部 \(a_i\)，因而包含 \(L\)。若 \(M/K\) 是 \(\Omega\) 中包含 \(L\) 的正规扩张，每个 \(m_{a_i,K}\) 都在 \(M\) 中分裂，所以 \(N_0\subseteq M\)。由\cref{b1:prop:normal-closure}的最小性，\(N=N_0\)。

记 \(M_0=K\bigl(\bigcup_{\sigma\in\operatorname{Hom}_K(L,\Omega)}\sigma(L)\bigr)\)。对每个 \(K\) 嵌入 \(\sigma:L\hookrightarrow\Omega\)，\(\sigma(a_i)\) 是 \(m_{a_i,K}\) 的根，故 \(\sigma(L)\subseteq N\)，进而 \(M_0\subseteq N\)。反过来，任取 \(m_{a_i,K}\) 在 \(\Omega\) 中的根 \(\beta\)。由\cref{b1:lem:simple-embedding-extension}，映射 \(a_i\mapsto\beta\) 给出 \(K(a_i)\hookrightarrow\Omega\) 的 \(K\) 嵌入；再由\cref{b1:thm:embedding-extension}将它延拓为 \(\sigma:L\hookrightarrow\Omega\)。于是 \(\beta\in\sigma(L)\subseteq M_0\)，所有这些最小多项式的根都属于 \(M_0\)，故 \(N\subseteq M_0\)。两边相等，而右边只取决于 \(L,K,\Omega\)，不取决于所选生成元。
\end{solution}

\begin{exercise}\label{b1:ex:15-14}
设 \(K\subseteq E\subseteq L\)，\(L/K\) 正规。证明：\(E/K\) 正规，当且仅当 \(L\) 的每个 \(K\) 自同构都将 \(E\) 映到自身。
\end{exercise}
\begin{solution}
若 \(E/K\) 正规，任意自同构的限制都是 \(E\) 到共同代数闭包的 \(K\) 嵌入，由\cref{b1:thm:normal-embedding-criterion}其像为 \(E\)。反过来，给定 \(E\) 到代数闭包的任意 \(K\) 嵌入，\cref{b1:cor:intermediate-embedding-extension}将它延拓到 \(L\) 的 \(K\) 自同构；按假设它把 \(E\) 映到自身。因此所有这样的嵌入都保持 \(E\)，再次用正规性判据得 \(E/K\) 正规。这里无需假定扩张可分。
\end{solution}
